%  06_experiments.tex
\section{Experiments and Results}
\label{sec:experiments}

\subsection{Leaky vs.\ honest evaluation}

We first quantify the impact of the evaluation protocol.  A Siamese network was trained on
the rousettus random-background dataset once under a \emph{random-pair} split (all
identities available in all partitions) and once under the \emph{identity-disjoint}
protocol.  Results are reported in Table~\ref{tab:leakage}.

\begin{table}[h]
\centering
\caption{Impact of evaluation protocol.  Both rows use the same Siamese architecture on
rousettus data; only the splitting strategy differs.  The random-pair row is from the
original (legacy) random-pair pipeline---the current codebase enforces identity-disjoint
splits by construction---while the identity-disjoint row is from the re-run on the fixed
code.}
\label{tab:leakage}
\begin{tabular}{lccc}
\toprule
Protocol   & ROC-AUC & F\textsubscript{1} & Recall \\
\midrule
Random-pair (leaky)        & 0.982 & 0.987 & 0.986 \\
Identity-disjoint (correct) & \bestval{0.858} & \bestval{0.707} & \bestval{0.734} \\
\bottomrule
\end{tabular}
\end{table}

The 12-point drop in ROC-AUC and 28-point drop in F\textsubscript{1} confirm that
random-pair splitting substantially overstates real-world performance.  All results
below use the identity-disjoint protocol exclusively.

\subsection{Model family comparison on rousettus}

Table~\ref{tab:rousettus} presents test-set results for the three model families on the
rousettus dataset under both the random-composited and green-screen background
conditions (3 unseen test identities in each).

\begin{table}[h]
\centering
\caption{Test results on \rousettus (identity-disjoint).  ArcFace v1 = paper-scale
default hyperparameters; ArcFace v2 / AdaFace = tuned hyperparameters.
ROC-AUC for \emph{Siamese} is computed from its native pairwise head; for
ArcFace/AdaFace it is cosine-similarity verification on the embeddings, so the two
are \textbf{not directly comparable} (Siamese cosine ROC-AUC, secondary: random
$0.946$, green $0.959$).  ``---'' = not applicable.  \bestval{Bold} marks the best
identification metric per background.  All rows except those marked $^{\dagger}$ are
significant at $p<0.05$ (permutation test, $n=100$).}
\label{tab:rousettus}
\setlength{\tabcolsep}{4pt}
\begin{tabular}{llccccc}
\toprule
Model & BG & ROC-AUC & F\textsubscript{1} & \topone & \mAP & TAR@\(10^{-3}\) \\
\midrule
Siamese & random & 0.858 & 0.707 & --- & --- & 0.012 \\
Siamese & green  & 0.669 & 0.797 & --- & --- & 0.001 \\
ArcFace (v2) & random & 0.665$^{\dagger}$ & --- & 0.532 & 0.731 & 0.074 \\
ArcFace (v2) & green  & 0.847 & --- & \bestval{0.846} & \bestval{0.915} & 0.067 \\
AdaFace & random & 0.434$^{\dagger}$ & --- & 0.274 & 0.560 & 0.003 \\
AdaFace & green  & 0.753 & --- & 0.556 & 0.762 & 0.006 \\
ArcFace (v1) & random & 0.500$^{\dagger}$ & --- & 0.591 & 0.757 & 0.000 \\
\bottomrule
\end{tabular}
\end{table}

On the \emph{random} background, the Siamese network leads decisively (pairwise-head
ROC-AUC $0.858$, $p<0.05$), while the tuned embedding models stay at or near chance and
are \emph{not} statistically significant (ArcFace $0.665$, $p=1.0$; AdaFace $0.434$,
$p=0.78$).  On the \emph{green-screen} background, however, the embedding models recover
sharply and become significant: ArcFace reaches ROC-AUC $0.847$ with Top-1 $0.846$ and
mAP $0.915$, exceeding the Siamese network on the identification metrics.  This
background dependence---rather than identity count alone---is the dominant factor in the
embedding models' behaviour (Section~\ref{sec:discussion}).  ArcFace v1 with paper-scale
hyperparameters still collapses to ROC-AUC $= 0.50$ (it degenerates to predicting every
pair ``same'', recall $1.0$), confirming that the large learning rate ($\eta = 0.1$)
destroys the ImageNet features before any face-specific knowledge is acquired.

\subsection{Model family comparison on mauritius}

Table~\ref{tab:mauritius} extends the comparison to the mauritius species across three
background conditions.

\begin{table}[h]
\centering
\caption{Test results on \mauritius by background condition (identity-disjoint).
``Siamese'' uses Focal loss; ``ArcFace''/``AdaFace'' use tuned angular losses.
As in Table~\ref{tab:rousettus}, Siamese ROC-AUC is pairwise-head and the embedding
ROC-AUC is cosine-based (not directly comparable; Siamese cosine ROC-AUC, secondary:
green $0.891$, original $0.927$, random $0.820$).  \bestval{Bold} marks the best
identification metric per background.  $^{\dagger}$ = not significant at $p<0.05$.}
\label{tab:mauritius}
\setlength{\tabcolsep}{4pt}
\begin{tabular}{llccccc}
\toprule
Model & BG & ROC-AUC & F\textsubscript{1} & \topone & \mAP & TAR@\(10^{-3}\) \\
\midrule
Siamese  & green    & 0.710 & 0.689 & --- & --- & 0.014 \\
Siamese  & original & 0.834 & 0.753 & --- & --- & 0.005 \\
Siamese  & random   & 0.806 & 0.686 & --- & --- & 0.007 \\
ArcFace  & green    & 0.843 & --- & 0.740 & 0.870 & 0.282 \\
AdaFace  & green    & 0.902 & --- & \bestval{0.916} & \bestval{0.955} & 0.383 \\
ArcFace  & original & 0.791 & --- & 0.576 & 0.774 & 0.145 \\
AdaFace  & original & 0.731 & --- & 0.634 & 0.785 & 0.151 \\
ArcFace  & random   & 0.578$^{\dagger}$ & --- & 0.479 & 0.697 & 0.021 \\
AdaFace  & random   & 0.539 & --- & 0.459 & 0.682 & 0.006 \\
\bottomrule
\end{tabular}
\end{table}

The Mauritius results sharpen the picture.  On the controlled \emph{green-screen}
condition the embedding models are now the strongest: AdaFace attains ROC-AUC $0.902$,
Top-1 $0.916$, and mAP $0.955$, surpassing both ArcFace and the Siamese network and
demonstrating that---with the corrected AdaFace margin, ImageNet normalisation, and
matched $112$-px evaluation---angular-margin embeddings generalise well to unseen
identities even at $\le 13$ training individuals \emph{when the background is coherent}.
On the \emph{original} background the embedding models remain solid (ArcFace $0.791$,
AdaFace $0.731$), whereas on \emph{random}-composited backgrounds they fall back toward
chance (ArcFace $0.578$, not significant; AdaFace $0.539$).  The Siamese network, by
contrast, is comparatively stable across all three backgrounds (pairwise-head ROC-AUC
$0.71$--$0.83$; cosine ROC-AUC $0.82$--$0.93$) and significant throughout.  Background
condition, not identity count, is thus the variable that most strongly separates the two
families.

\subsection{ArcFace hyperparameter sensitivity}

Figure~\ref{fig:arcface_collapse} illustrates the training dynamics under the two ArcFace
configurations.

\begin{figure}[t]
\centering
\begin{tikzpicture}
\begin{axis}[
  name=collapseA,
  batplot,
  height=5cm, width=0.46\linewidth,
  title={\small ArcFace v1 (collapsed, SGD $\eta{=}0.1$)},
  xlabel={Epoch},
  ylabel={Metric},
  ymin=0, ymax=1,
  xtick={1,4,8},
  legend pos=south east
]
\addplot[thick, colArcFace, mark=*, mark size=1.5]
  coordinates {(1,0.126)(2,0.346)(3,0.654)(4,0.815)(5,0.819)(6,0.912)(7,0.924)(8,0.928)};
\addlegendentry{train acc}
\addplot[thick, colSiamese, dashed, mark=square*, mark size=1.5]
  coordinates {(1,0.50)(2,0.50)(3,0.50)(4,0.50)(5,0.50)(6,0.50)(7,0.351)(8,0.50)};
\addlegendentry{val ROC-AUC}
\end{axis}

\begin{axis}[
  at={(collapseA.outer north east)}, anchor=outer north west,
  xshift=1.0cm,
  batplot,
  height=5cm, width=0.46\linewidth,
  title={\small ArcFace v2 (tuned, Adam $\eta{=}10^{-3}$)},
  xlabel={Epoch},
  ymin=0, ymax=1,
  xtick={1,5,10,13},
  legend pos=south east
]
\addplot[thick, colArcFace, mark=*, mark size=1.5]
  coordinates {(1,0.451)(2,0.949)(3,0.957)(4,0.979)(5,0.984)(6,0.992)(7,0.986)
               (8,0.969)(9,0.965)(10,0.988)(11,0.986)(12,0.959)(13,0.969)};
\addlegendentry{train acc}
\addplot[thick, colSiamese, dashed, mark=square*, mark size=1.5]
  coordinates {(1,0.579)(2,0.389)(3,0.512)(4,0.599)(5,0.648)(6,0.665)(7,0.652)
               (8,0.634)(9,0.419)(10,0.476)(11,0.661)(12,0.545)(13,0.566)};
\addlegendentry{val ROC-AUC}
\end{axis}
\end{tikzpicture}
\caption{Training dynamics of ArcFace on rousettus (random background).
\emph{Left}: the default (paper-scale) SGD configuration fails---train accuracy climbs
to ${\sim}0.93$ while val ROC-AUC stays pinned at chance ($0.50$), the hallmark of a
closed-set classifier that fits the training identities but learns no transferable
verification geometry.
\emph{Right}: the tuned Adam configuration reaches near-perfect train accuracy but val
ROC-AUC oscillates and peaks only around $0.66$, revealing the generalisation gap to
unseen test identities under random-background compositing.}
\label{fig:arcface_collapse}
\end{figure}

\subsection{Learning curves}

Figure~\ref{fig:learning_curves} shows per-epoch validation curves for all three models on
rousettus (random background).

\begin{figure}[t]
\centering
\begin{tikzpicture}
\begin{axis}[
  batplot,
  height=6cm,
  width=0.92\linewidth,
  xlabel={Epoch},
  ylabel={Val ROC-AUC},
  ymin=0.3, ymax=1.0,
  legend pos=south east,
  title={\small Validation ROC-AUC during training --- Rousettus (random background)}
]

% Siamese (10 epochs)
\addplot[thick, colSiamese, mark=*, mark size=2]
  coordinates {
    (1,0.811)(2,0.855)(3,0.853)(4,0.859)(5,0.861)
    (6,0.858)(7,0.862)(8,0.863)(9,0.847)(10,0.838)
  };
\addlegendentry{Siamese (Focal)}

% ArcFace v2 (13 epochs)
\addplot[thick, colArcFace, dashed, mark=square*, mark size=2]
  coordinates {
    (1,0.579)(2,0.389)(3,0.512)(4,0.599)(5,0.648)(6,0.665)(7,0.652)
    (8,0.634)(9,0.419)(10,0.476)(11,0.661)(12,0.545)(13,0.566)
  };
\addlegendentry{ArcFace (tuned)}

% AdaFace (8 epochs)
\addplot[thick, colAdaFace, dotted, mark=triangle*, mark size=2]
  coordinates {
    (1,0.641)(2,0.408)(3,0.424)(4,0.450)(5,0.415)(6,0.419)(7,0.477)(8,0.434)
  };
\addlegendentry{AdaFace (tuned)}

\end{axis}
\end{tikzpicture}
\caption{Validation ROC-AUC per training epoch on rousettus (random background).  The
Siamese network converges quickly and stably to ${\sim}0.86$.  ArcFace and AdaFace show
high epoch-to-epoch variance and never stabilise on this random-background condition,
with AdaFace decaying toward chance, reflecting the difficulty of learning transferable
angular geometry under random-background compositing.  Early stopping terminates the
embedding runs after $8$--$13$ epochs (best-validation checkpoint restored for test).}
\label{fig:learning_curves}
\end{figure}

\subsection{Family and background comparison (summary bar chart)}

Figure~\ref{fig:bar_summary} provides a visual summary of ROC-AUC across all experimental
conditions.

\begin{figure}[t]
\centering
\begin{tikzpicture}
% --- LEFT: Rousettus ---
\begin{axis}[
  name=ax1,
  ybar, bar width=11pt,
  height=5.5cm, width=0.44\linewidth,
  title={\small Rousettus (random BG)},
  ylabel={Test ROC-AUC},
  symbolic x coords={Siamese,ArcFace,AdaFace},
  xtick=data,
  ymin=0.4, ymax=0.95,
  ymajorgrids=true,
  grid style={line width=0.3pt, draw=gray!30},
  tick label style={font=\small},
  label style={font=\small},
  nodes near coords,
  nodes near coords style={font=\tiny,
    /pgf/number format/fixed,
    /pgf/number format/precision=3}
]
\addplot[fill=colSiamese!70, draw=colSiamese]
  coordinates {(Siamese,0.858) (ArcFace,0.665) (AdaFace,0.434)};
\end{axis}

% --- RIGHT: Mauritius ---
\begin{axis}[
  at={(ax1.outer north east)}, anchor=outer north west,
  xshift=1.2cm,
  ybar, bar width=8pt,
  height=5.5cm, width=0.52\linewidth,
  title={\small Mauritius (all conditions)},
  xlabel={Model / Condition},
  symbolic x coords={Siam-grn,Siam-ori,Siam-rnd,Arc-grn,Ada-grn,Arc-rnd},
  xtick=data,
  xticklabel style={font=\tiny, rotate=30, anchor=north east},
  ymin=0.4, ymax=1.05,
  ymajorgrids=true,
  grid style={line width=0.3pt, draw=gray!30},
  tick label style={font=\small},
  label style={font=\small},
  nodes near coords,
  nodes near coords style={font=\tiny,
    /pgf/number format/fixed,
    /pgf/number format/precision=3}
]
\addplot[fill=colSiamese!70, draw=colSiamese]
  coordinates {
    (Siam-grn,0.710) (Siam-ori,0.834) (Siam-rnd,0.806)
    (Arc-grn,0.843) (Ada-grn,0.902) (Arc-rnd,0.578)
  };
\end{axis}

\end{tikzpicture}
\caption{Test ROC-AUC for the major conditions.
\emph{Left}: rousettus, random background.
\emph{Right}: mauritius across background conditions and model families.
``Siam'' = Siamese, ``Arc'' = ArcFace, ``Ada'' = AdaFace; suffixes grn/ori/rnd indicate
green/original/random backgrounds.  Siamese bars use the pairwise-head ROC-AUC and the
ArcFace/AdaFace bars use cosine-verification ROC-AUC, so heights are not directly
comparable across families; the figure is intended to show the within-family
background trend, in particular the embedding models' recovery on green/original and
their collapse on random.}
\label{fig:bar_summary}
\end{figure}

\subsection{Reproducibility}

Table~\ref{tab:repro} documents the effect of the DataLoader seeding fix on run-to-run
reproducibility for ArcFace.  Without an explicit \texttt{torch.Generator} in the
\texttt{DataLoader}, the global RNG state (already advanced by model initialisation and
EMA setup) causes non-deterministic batch ordering even with \texttt{deterministic=True}.

\begin{table}[h]
\centering
\caption{Run-to-run reproducibility for ArcFace (rousettus, random BG).
Repro-1/2: before DataLoader seed fix; Repro-3/4: after fix.}
\label{tab:repro}
\begin{tabular}{lcc}
\toprule
Run & test ROC-AUC & test Top-1 \\
\midrule
repro-1 (pre-fix)   & 0.528 & 0.326 \\
repro-2 (pre-fix)   & 0.563 & 0.405 \\
repro-3 (post-fix)  & 0.552 & 0.361 \\
repro-4 (post-fix)  & 0.552 & 0.361 \\
\bottomrule
\end{tabular}
\end{table}

Post-fix, the two identical runs (repro-3 and repro-4) produce bit-exact identical metrics,
confirming that the fix resolves non-determinism.
